\chapter{Оригинальная часть работы}\label{ch:original}

\section{Решение обратной спектроскопической задачи для молекулы германа}\label{sec:ch2/Ge}

\section{Решение обратной спектроскопической задачи для молекулы силана}\label{sec:ch2/Si}

В данном параграфе рассматривается молекула силана SiH$_4$, то есть третий, после метана и германа,
представитель семейства молекул XY$_4$ с тетраэдрической симметрией T$_d$, для которого была решена
обратная спектроскопическая задача в чисто колебательном приближении. Под обратной задачей здесь
понимается восстановление набора эффективных спектроскопических параметров по совокупности
экспериментально определенных центров полос и последующее использование этих параметров для
расчета энергетической структуры изотопологов ${}^{M}$SiH$_4$ ($M=28,29,30$) в области до
9000 см$^{-1}$.

Актуальность такой постановки задачи связана как с фундаментальными, так и с прикладными
причинами. Силан и его изотопологи используются при моделировании атмосфер газовых гигантов и их
спутников, прежде всего Юпитера, Сатурна и Титана \cite{003,004,005,006,007,008,009}; линии SiH$_4$
также важны при интерпретации спектров околозвездных оболочек, в частности объекта IRC+10216
\cite{010,011,012}. Кроме того, контроль содержания силана необходим в технологических процессах
получения высокочистого кремния и тонких кремниевых пленок \cite{001,002}. Поэтому для таких задач
требуются не только качественные сведения о спектре, но и надежные численные значения центров полос,
ангармонических параметров, параметров резонансного взаимодействия и параметров тетраэдрического
расщепления.

В современных исследованиях параметры поверхности потенциальной энергии могут быть получены двумя
основными путями. Первый путь состоит в прямом квантово-химическом, или <<ab~initio>>, расчете
\cite{abinitio_1,abinitio_3,abinitio_5,abinitio_6,abinitio_8,abinitio_9,abinitio_10,abinitio_11}. Несмотря на
значительный прогресс вычислительных методов, для задач спектроскопии высокого разрешения точность
таких расчетов, как правило, остается ниже точности полуэмпирического описания, основанного на
экспериментальных данных. Второй путь заключается в определении параметров эффективного
гамильтониана по измеренным спектральным характеристикам \cite{kuchitsu,nielsen,locmode,MWour}. Именно
этот подход применялся в настоящей работе: из экспериментальных центров полос восстанавливались
частоты гармонического приближения, ангармонические коэффициенты, параметры тетраэдрического
расщепления и параметры резонансных взаимодействий.

\subsection{Модель эффективного колебательного гамильтониана}

Для молекул типа асимметричного и симметричного волчка аналитические формы базисных функций и
матричных элементов гамильтониана хорошо известны \cite{0101,0102,0103,0104,8855,0105}. Для молекул
типа сферического волчка аналогичная задача существенно сложнее, поскольку необходимо учитывать
неприводимые представления группы T$_d$ и корректно строить симметризованный колебательный базис.
В работе \cite{hecht} были получены основные соотношения для низших состояний молекул XY$_4$, а в
\cite{ulen1,ulen2} были построены аналитические выражения для симметризованных колебательных функций
и матричных элементов, необходимые для анализа полиад вплоть до $N=4$.

В используемом подходе симметризованная колебательная функция молекулы XY$_4$ записывается в виде
произведения функции невырожденного осциллятора, функции дважды вырожденного осциллятора и
симметризованной комбинации функций двух трижды вырожденных осцилляторов:
\begin{equation}
\left| v_1; v_2 l_2 \Gamma_2; v_3 l_3, v_4 l_4, l n_l \Gamma_{34}; n\Gamma\sigma \right\rangle =
\left| v_1 \right\rangle
\left(\left| v_2 l_2 \Gamma_2 \right\rangle \otimes
\left| v_3 l_3, v_4 l_4, l n_l \Gamma_{34} \right\rangle\right)^{n\Gamma}_{\sigma} .
\label{eq:si_basis_function}
\end{equation}
Здесь $v_i$ обозначают колебательные квантовые числа, $l_i$ --- соответствующие колебательные
моменты, $\Gamma$ --- тип симметрии состояния, а индекс $\sigma$ задает строку неприводимого
представления. При построении базиса используются методы теории неприводимых тензорных наборов
\cite{Fano1959,Wigner1965,Varshalovich1988}. Для части функции, связанной с двумя трижды
вырожденными координатами, существенную роль играют коэффициенты редукции
${}^{l(a)}G^m_{n_l\Gamma_{34}\sigma_{34}}$, которые переводят функции группы $O(3)$ к функциям
точечной группы T$_d$. Первоначально такие коэффициенты использовались в табличной форме
\cite{moret,mor2,Champ1,Champ2}; аналитическая форма их записи была предложена в \cite{chegl} и затем
использована при построении полного набора функций в \cite{ulen1,ulen2}.

Для заданного набора симметризованных функций матричный элемент эффективного колебательного
гамильтониана можно представить как
\begin{equation}
H_{vv'}^{n\Gamma,n'\Gamma'} = \widetilde{E}_{v}\,\delta_{vv'}\delta_{nn'}\delta_{\Gamma\Gamma'}
+ W_{vv'}^{n\Gamma,n'\Gamma'} .
\label{eq:si_hamiltonian_matrix}
\end{equation}
Первое слагаемое задает невозмущенную ангармоническую энергию состояния $v\equiv(v_1v_2v_3v_4)$,
которая в общем виде записывается как
\begin{equation}
\begin{aligned}
\widetilde{E}_{v} ={}&
\sum_i \omega_i\left(v_i+\frac{d_i}{2}\right) \\
&+\sum_{i,j\geq i} x_{ij}
\left(v_i+\frac{d_i}{2}\right)
\left(v_j+\frac{d_j}{2}\right) \\
&+\sum_{i,j\geq i,k\geq j} y_{ijk}
\left(v_i+\frac{d_i}{2}\right)
\left(v_j+\frac{d_j}{2}\right)
\left(v_k+\frac{d_k}{2}\right)
+\ldots .
\end{aligned}
\label{eq:si_unperturbed_energy}
\end{equation}
Здесь $\omega_i$ --- гармонические частоты, $x_{ij}$ и $y_{ijk}$ --- ангармонические параметры,
$d_i$ равно 1, 2 или 3 для невырожденного, дважды и трижды вырожденного колебания соответственно.
Второе слагаемое $W_{vv'}^{n\Gamma,n'\Gamma'}$ описывает тетраэдрическое расщепление при $v=v'$ и
резонансные взаимодействия при $v\neq v'$. Типичные выражения для матричных элементов $W_{vv'}^{n\Gamma,n'\Gamma'}$,
использованные при построении эффективного колебательного гамильтониана,
приведены в табл.~\ref{tab:si_hamiltonian_elements_1}
и табл.~\ref{tab:si_hamiltonian_elements_2}.

% Floating table for part2; requires graphicx.
\begin{table}[p]
\centering
\caption{Тетраэдрическое расщепление для полосы $3\nu_2+3\nu_4(F_1)$.}
\label{tab:si_hamiltonian_elements_1}
\begingroup
%%%%%%%%%%%%%%%%%%%%%%%%%%%%%%%%%%%%%%%%%%%%%%%%%%%%%%%%%%%%
\def\A{$9G_{22}+12G_{44}-4T_{44}$ }
\def\ab{$0$ }
\def\ac{$0$ }
\def\ad{$-\frac{24}{5}\sqrt{10}T_{24}$ }
\def\ai{$0$ }
\def\af{$-\frac{24}{5}\sqrt{15}T_{24}$ }
\def\B{$9G_{22}+2G_{44}$ }
\def\bc{$8\sqrt{6}T_{44}$}
\def\bd{$\frac{72}{5}\sqrt{3}T_{24}$}
\def\be{$-\frac{24}{5}\sqrt{10}T_{24}$}
  \def\br{$-\frac{72}{5}\sqrt{2}T_{24}$ }
\def\C{$9G_{22}+12G_{44}+12T_{44}$}
\def\cd{$\frac{72}{5}\sqrt{2}T_{24}$}
\def\ce{$\frac{24}{5}\sqrt{15}T_{24}$}
\def\cf{$-\frac{48}{5}\sqrt{3}T_{24}$ }
\def\D{$G_{22}+2G_{44}-\frac{144}{5}T_{24}$}
\def\de{$-\frac{16}{5}\sqrt{30}T_{24}$}
\def\df{$8\sqrt{6}(\frac{6}{5}T_{24}-T_{44})$ }
\def\E{$G_{22}+12G_{44}-4T_{44}$}
\def\ef{$-\frac{48}{5}\sqrt{5}T_{24}$ }
\def\F{$G_{22}+12G_{44}$}
\def\FF{$-\frac{96}{5}T_{24}+12T_{44}$}
  \def\1{$l_2=3,l_4=3$ }
  \def\2{$l_2=3,l_4=1$ }
  \def\3{$l_2=3,l_4=3$ }
  \def\4{$l_2=1,l_4=1$ }
  \def\5{$l_2=1,l_4=3$ }
  \def\6{$l_2=1,l_4=3$ }
\scriptsize
\rotatebox{90}{%
\resizebox{0.83\textheight}{!}{%
\begin{tabular}{|r c|c|c|c|c|c l|}\hline
&\multicolumn{6}{c}{} & \\
&\multicolumn{6}{c}{$3\nu_2+3\nu_4(F_1)$} & \\
&\multicolumn{6}{c}{} & \\
\hline
&&&&&&&\\
% A       B      C      D      E      F
&\1     & \2    & \3   & \4   & \5   & \6  & \\
&&&&&&&\\
\hline
&&&&&&&\\
& \A    & \ab   & \ac  & \ad  & \ai  & \af  & \\
&&&&&&&\\
\hline
&&&&&&&\\
&       & \B    & \bc  & \bd  & \be  & \br  & \\
&&&&&&&\\
\hline
&&&&&&&\\
&       &       & \C   & \cd  & \ce  & \cf &  \\
&&&&&&&\\
\hline
&&&&&&&\\
&       &       &      & \D   & \de  & \df  & \\
&&&&&&&\\
\hline
&&&&&&&\\
&       &       &      &      & \E   & \ef  & \\
&&&&&&&\\
\hline
&&&&&&&\\
&       &       &      &      &      & \F   & \\
&       &       &      &      &      & \FF  & \\
&&&&&&&\\
\hline
\end{tabular}
}%
}
\endgroup
\end{table}


% Floating table for part2; requires graphicx.
\begin{table}[p]
\centering
\caption{Тетраэдрическое расщепление и резонансные взаимодействия для полос $\nu_2+2\nu_3+\nu_4$ и $\nu_1+\nu_2+\nu_3+\nu_4(F_1)$.}
\label{tab:si_hamiltonian_elements_2}
\begingroup
  \def\A{$G_{22}+6G_{33}+2G_{44}$ }
  \def\AR{$-8T_{23}-8T_{24}+12T_{33}$ }
  \def\AAA{$+\frac{8}{3}S_{34}+4T_{34}$ }
  \def\ab{$\frac{1}{5}\sqrt{10}(8T_{23}+8T_{24}$ }
  \def\aab{$-12T_{33}-2T_{34}$ }
  \def\aaab{$+3G_{34}-\frac{49}{12}S_{34})$ }
  \def\ac{$-\frac{8}{3}i\sqrt{6}T_{23}$ }
  \def\ad{$0$ }
  \def\ai{$\frac{1}{5}\sqrt{15}(+\frac{13}{4}S_{34}-8T_{24}$ }
    \def\aae{$+12T_{33}+2G_{34}$ }
    \def\aaae{$+\frac{16}{3}T_{34}-8T_{23})$ }
  \def\af{$0$ }
  \def\ag{$0$ }
  \def\B{$G_{22}+6G_{33}+2G_{44}$ }
  \def\BB{$-\frac{28}{5}T_{23}-3G_{34}$ }
  \def\BBB{$-\frac{4}{5}T_{24}+\frac{14}{3}S_{34}$ }
  \def\bc{$\frac{16}{15}i\sqrt{15}T_{23}$}
  \def\bd{$-\frac{4}{5}\sqrt{30}(T_{23}+T_{24})$}
  \def\be{$\frac{4}{5}\sqrt{6}(T_{23}+3T_{24}$}
  \def\bbe{$-5T_{33}-\frac{5}{6}T_{34})$}
  \def\br{$0$ }
  \def\bg{$-\frac{1}{20}\sqrt{15}K_{1333}$ }
  \def\C{$G_{22}+6G_{33}+2G_{44}$}
  \def\CC{$-\frac{16}{3}T_{23}+\frac{16}{3}T_{24}+8T_{33}$}
  \def\CCC{$-\frac{17}{3}S_{34}-\frac{8}{3}T_{34}$}
  \def\cd{$\frac{1}{3}i\sqrt{2}(4T_{34}-S_{34}-2G_{34})$}
  \def\ce{$\frac{8}{5}i\sqrt{10}T_{23}$}
  \def\cf{$0$ }
  \def\cg{$0$ }
  \def\D{$G_{22}+6G_{33}+2G_{44}$}
  \def\DD{$+\frac{8}{3}(T_{23}+T_{34}-T_{24}$}
  \def\DDD{$-2T_{33})-\frac{13}{3}S_{34}+\frac{2}{3}G_{34}$}
  \def\de{$-\frac{8}{5}\sqrt{5}(T_{23}+T_{24})$}
  \def\df{$-\frac{1}{12}i\sqrt{6}K_{1333}$ }
  \def\dg{$0$ }
  \def\E{$G_{22}+6G_{33}+2G_{44}$}
  \def\EE{$-\frac{32}{5}T_{23}-\frac{16}{5}T_{24}+4T_{33}$}
  \def\EEE{$+\frac{4}{3}S_{34}+4T_{34}+2G_{34}$}
  \def\ef{$0$ }
  \def\eg{$-\frac{1}{20}\sqrt{10}K_{1333}$ }
  \def\F{$G_{22}+2G_{33}+2G_{44}$}
  \def\FF{$-4T_{23}-4T_{24}$}
  \def\FFF{$-\frac{10}{3}S_{34}-G_{34}$}
  \def\fg{$4i\sqrt{3}(T_{23}-T_{24})$ }
  \def\G{$G_{22}+2G_{33}+2G_{44}$}
  \def\GG{$+4T_{23}+4T_{24}$}
  \def\GGG{$+\frac{2}{3}S_{34}-4T_{34}+G_{34}$}
  \def\1{$l_2=1,l_3=2,l_4=1$ }
  \def\2{$l_2=1,l_3=2,l_4=1$ }
  \def\3{$l_2=1,l_3=2,l_4=1$ }
  \def\4{$l_2=1,l_3=2,l_4=1$ }
  \def\5{$l_2=1,l_3=2,l_4=1$ }
  \def\6{$l_2=1,l_3=1,l_4=1$ }
  \def\7{$l_2=1,l_3=1,l_4=1$ }
\tiny
\rotatebox{90}{%
\resizebox{0.92\textheight}{!}{%
\begin{tabular}{|r c|c|c|c|c l|r c|c l|}\hline
& \multicolumn{5}{c}{}                     &&&   \multicolumn{2}{c}{}                          &  \\
& \multicolumn{5}{c}{$\nu_2+2\nu_3+\nu_4(F_1)$} &&&   \multicolumn{2}{c}{$\nu_1+\nu_2+\nu_3+\nu_4(F_1)$} &  \\
& \multicolumn{5}{c}{}                     &&&   \multicolumn{2}{c}{}                          &  \\
\hline
&&&&&&&&&&\\
% A       B      C      D      E      F      G
&\1     & \2    & \3   & \4   & \5   &&& \6   & \7  & \\
&&&&&&&&&&\\
\hline
&&&&&&&&&&\\
& \A    & \ab   & \ac  & \ad  & \ai  &&& \af  & \ag & \\
& \AR   & \aab  &      &      & \aae &&&      &     & \\
& \AAA  & \aaab &      &      & \aaae&&&      &     & \\
&&&&&&&&&&\\
\hline
&&&&&&&&&&\\
&       & \B    & \bc  & \bd  & \be  &&& \br  & \bg  & \\
&       & \BB   &      &      & \bbe &&&      &      & \\
&       & \BBB  &      &      &      &&&      &      & \\
&&&&&&&&&&\\
\hline
&&&&&&&&&&\\
&       &       & \C   & \cd  & \ce  &&& \cf  & \cg & \\
&       &       & \CC  &      &      &&&      &     & \\
&       &       & \CCC &      &      &&&      &     & \\
&&&&&&&&&&\\
\hline
&&&&&&&&&&\\
&       &       &      & \D   & \de  &&& \df  & \dg & \\
&       &       &      & \DD  &      &&&      &     & \\
&       &       &      & \DDD &      &&&      &     & \\
&&&&&&&&&&\\
\hline
&&&&&&&&&&\\
&       &       &      &      & \E   &&& \ef  & \eg & \\
&       &       &      &      & \EE  &&&      &     & \\
&       &       &      &      & \EEE &&&      &     & \\
&&&&&&&&&&\\
\hline
&&&&&&&&&&\\
&       &       &      &      &      &&& \F   &     & \\
&       &       &      &      &      &&& \FF  & \fg & \\
&       &       &      &      &      &&& \FFF &     & \\
&&&&&&&&&&\\
\hline
&&&&&&&&&&\\
&       &       &      &      &      &&&      & \G  & \\
&       &       &      &      &      &&&      & \GG & \\
&       &       &      &      &      &&&      & \GGG& \\
&&&&&&&&&&\\
\hline
\end{tabular}
}%
}
\endgroup
\end{table}

Дополнительные матричные элементы, необходимые для описания тетраэдрического
расщепления и резонансных взаимодействий в высоколежащих полиадах,
представлены в табл.~\ref{tab:si_hamiltonian_elements_a1}
и табл.~\ref{tab:si_hamiltonian_elements_f2}.

% Non-floating appendix table; requires graphicx.
% Floating table for part2; requires graphicx.
\begin{table}[p]
\centering
\caption{Тетраэдрическое расщепление и резонансные взаимодействия для $A_1$-подуровней состояний $4\nu_2+\nu_3+\nu_4$, $4\nu_1+2\nu_2+3\nu_4$, $2\nu_2+\nu_3+3\nu_4$, $\nu_3+5\nu_4$ и $\nu_2+5\nu_4$.}
\label{tab:si_hamiltonian_elements_a1}
\begingroup
  \def\A{$16G_{22}+2G_{33}$ }
  \def\AR{$+2G_{44}+\frac{2}{3}S_{34}$ }
  \def\AAA{$+6T_{34}+G_{34}$ }
  \def\ab{$-16(T_{23}+T_{24})$ }
  \def\ac{$0$ }
  \def\ad{$0$ }
  \def\ai{$\frac{41}{32}\sqrt{\frac{3}{5}}d_{24t}$ }
  \def\af{$-\sqrt{\frac{2}{5}}d_{24t}$ }
  \def\ag{$0$ }
  \def\ah{$0$ }
  \def\aii{$-\frac{41}{32}\sqrt{\frac{1}{3}}d_{24t}$ }
  \def\aj{$0$ }
  \def\ak{$0$ }
  \def\al{$0$ }
  \def\am{$0$ }
  \def\an{$0$ }
  \def\B{$4G_{22}+2G_{33}$ }
  \def\BB{$+2G_{44}+\frac{2}{3}S_{34}$ }
  \def\BBB{$+6T_{34}+G_{34}$ }
  \def\bb{$16\sqrt{3}(T_{23}+T_{24})$ }
  \def\bc{$0$ }
  \def\bd{$0$ }
  \def\be{$\sqrt{\frac{15}{2}}d_{24s}$ }
  \def\br{$-\frac{3}{2}\sqrt{\frac{3}{5}}d_{24t}$ }
  \def\bg{$\sqrt{\frac{2}{5}}d_{24t}$ }
  \def\bh{$0$ }
  \def\bi{$0$ }
  \def\bj{$0$ }
  \def\bk{$0$ }
  \def\bl{$0$ }
  \def\bm{$0$ }
  \def\bn{$0$ }
  \def\C{$2G_{33}+2G_{44}$ }
  \def\CC{$+\frac{20}{3}S_{34}+2G_{34}$ }
  \def\cd{$0$ }
  \def\ce{$-\frac{23}{15}\sqrt{\frac{1}{5}}d_{24t}$ }
  \def\cf{$\sqrt{\frac{2}{15}}d_{24t}$ }
  \def\cg{$0$ }
  \def\ch{$\sqrt{10}d_{24s}$ }
  \def\ci{$-\frac{25}{48}d_{24t}$ }
  \def\cj{$0$ }
  \def\ck{$0$ }
  \def\cl{$0$ }
  \def\cm{$0$ }
  \def\cn{$0$ }
  \def\D{$12G_{44}+24T_{44}$ }
  \def\de{$0$ }
  \def\df{$0$ }
  \def\dg{$0$ }
  \def\dr{$0$ }
  \def\di{$0$ }
  \def\djr{$0$ }
  \def\dk{$0$ }
  \def\dl{$0$ }
  \def\dm{$0$ }
  \def\dn{$\frac{3}{\sqrt{2}}d_{24s}$ }
  \def\E{$4G_{22}+2G_{33}$ }
  \def\EE{$+12G_{44}+\frac{159}{16}T_{44}$ }
  \def\EEE{$+\frac{281}{160}S_{34}+\frac{489}{160}T_{34}$ }
  \def\EEEE{$-\frac{51}{32}G_{34}$ }
  \def\ef{$-\frac{1}{20}\sqrt{\frac{3}{2}}(290T_{44}$ }
  \def\eff{$+53T_{34}+23S_{34})$ }
  \def\eg{$\frac{29}{2}T_{23}+\frac{159}{10}T_{24}$ }
  \def\eh{$-\frac{46}{5}\sqrt{6}T_{24}$ }
  \def\ei{$\frac{1}{32}\sqrt{5}\left(6T_{44}+\frac{127}{20}S_{34}\right.$ }
  \def\eii{$\left.-\frac{67}{5}T_{34}-43G_{34}\right)$ }
  \def\ej{$-\frac{295}{64}\sqrt{\frac{1}{21}}d_{24t}$ }
  \def\ek{$-\frac{1}{64}\sqrt{15}d_{24t}$ }
  \def\el{$\frac{53}{80}\sqrt{\frac{1}{6}}d_{24t}$ }
  \def\er{$-\frac{23}{20}\sqrt{\frac{1}{14}}d_{24t}$ }
  \def\en{$0$ }
  \def\F{$4G_{22}+2G_{33}$ }
  \def\FF{$+2G_{44}+\frac{6}{5}S_{34}$ }
  \def\FFF{$+\frac{54}{5}T_{34}+G_{34}$ }
  \def\fg{$-\frac{48}{5}\sqrt{6}T_{24}$ }
  \def\fh{$16\left(T_{23}+\frac{9}{5}T_{24}\right)$ }
  \def\fr{$\sqrt{30}\frac{1}{24}(17T_{34}-6T_{44}-5S_{34})$ }
  \def\fj{$0$ }
  \def\fk{$0$ }
  \def\fl{$-\frac{9}{10}d_{24t}$ }
  \def\fm{$\frac{2}{5}\sqrt{\frac{7}{3}}d_{24t}$ }
  \def\fn{$0$ }
  \def\G{$2G_{33}+12G_{44}$ }
  \def\GG{$+12T_{44}+2S_{34}$ }
  \def\GGG{$+6T_{34}+3G_{34}$ }
  \def\gh{$-4\sqrt{6}(2T_{44}+T_{34})$ }
  \def\gi{$\frac{1}{2}\sqrt{5}(T_{23}-17T_{24})$ }
  \def\gj{$0$ }
  \def\gk{$0$ }
  \def\gl{$\frac{3}{\sqrt{2}}d_{24s}$ }
  \def\gm{$0$ }
  \def\gn{$0$ }
  \def\HH{$2G_{33}+2G_{44}$ }
  \def\HHH{$+12S_{34}-2G_{34}$ }
  \def\hi{$-10\sqrt{\frac{10}{3}}T_{24}$ }
  \def\hj{$0$ }
  \def\hk{$0$ }
  \def\hl{$0$ }
  \def\hm{$\sqrt{7}d_{24s}$ }
  \def\hn{$0$ }
  \def\I{$4G_{22}+2G_{33}$ }
  \def\II{$+12G_{44}-\frac{167}{48}T_{44}$ }
  \def\III{$+\frac{65}{32}S_{34}-\frac{95}{32}T_{34}$ }
  \def\IIII{$+\frac{1}{32}G_{34}$ }
  \def\ir{$-\frac{121}{192}\sqrt{\frac{5}{21}}d_{24t}$ }
  \def\ik{$-\frac{91}{64}\sqrt{\frac{1}{3}}d_{24t}$ }
  \def\il{$-\frac{17}{48}\sqrt{\frac{5}{6}}d_{24t}$ }
  \def\im{$-\sqrt{\frac{5}{14}}\frac{5}{12}d_{24t}$ }
  \def\io{$0$ }
  \def\J{$2G_{33}+30G_{44}$ }
  \def\JJ{$+20T_{44}+\frac{95}{18}S_{34}$ }
  \def\JJJ{$+\frac{35}{6}T_{34}-\frac{15}{4}G_{34}$ }
  \def\jk{$\sqrt{35}\left(\frac{20}{7}T_{44}+\frac{1}{24}S_{34}\right.$ }
  \def\jkk{$\left.-\frac{1}{6}T_{34}-\frac{3}{4}G_{34}\right)$ }
  \def\jl{$-\sqrt{14}\left(\frac{40}{7}T_{44}+\frac{5}{9}S_{34}-\frac{20}{21}T_{34}\right)$ }
  \def\jm{$40\sqrt{\frac{2}{3}}T_{44}$ }
  \def\jn{$0$ }
  \def\K{$2G_{33}+30G_{44}$ }
  \def\KK{$-20T_{44}+\frac{23}{6}S_{34}$ }
  \def\KKK{$-\frac{21}{2}T_{34}+\frac{11}{4}G_{34}$ }
  \def\kl{$\sqrt{10}\left(8T_{44}-\frac{1}{3}S_{34}+\frac{4}{3}T_{34}\right)$ }
  \def\km{$8\sqrt{\frac{30}{7}}T_{44}$ }
  \def\kn{$0$ }
  \def\LL{$2G_{33}+12G_{44}$ }
  \def\LLL{$+28T_{44}+\frac{26}{9}S_{34}$ }
  \def\LLLL{$+\frac{26}{3}T_{34}+3G_{34}$ }
  \def\lm{$-\sqrt{21}\left(\frac{208}{21}T_{44}+\frac{24}{7}T_{34}\right)$ }
  \def\lr{$0$ }
  \def\M{$2G_{33}+2G_{44}$ }
  \def\MM{$+\frac{52}{3}S_{34}-2G_{34}$ }
  \def\mn{$0$ }
  \def\N{$12G_{44}-56T_{44}$}
  \def\1{$l_2=4,l_3=1,l_4=1$ }%
  \def\2{$l_2=2,l_3=1,l_4=1$ }%
  \def\3{$l_2=0,l_3=1,l_4=1$ }%
  \def\4{$l_2=0,l_3=0,l_4=3$ }%
  \def\5{$l_2=2,l_3=1,l_4=3$ }%
  \def\6{$l_2=2,l_3=1,l_4=1$ }%
  \def\7{$l_2=0,l_3=1,l_4=3$ }%
  \def\8{$l_2=0,l_3=1,l_4=1$ }%
  \def\9{$l_2=2,l_3=1,l_4=3$ }%
  \def\la{$l_2=0,l_3=1,l_4=5$ }%
  \def\lb{$l_2=0,l_3=1,l_4=5$ }%
  \def\lc{$l_2=0,l_3=1,l_4=3$ }%
  \def\ld{$l_2=0,l_3=1,l_4=1$ }%
  \def\li{$l_2=0,l_3=0,l_4=3$ }%
\tiny
\rotatebox{90}{%
\resizebox{0.80\textheight}{!}{%
\begin{tabular}{|r c|c|c l|r c l|r c|c|c|c|c l|r c|c|c|c l|r c l|}\hline
& \multicolumn{3}{c}{}                     &&&   \multicolumn{1}{c}{}                          &&&   \multicolumn{5}{c}{}                          &&&   \multicolumn{4}{c}{}                          &&&   \multicolumn{1}{c}{}                          &  \\
& \multicolumn{3}{c}{$4\nu_2+\nu_3+\nu_4$} &&&   \multicolumn{1}{c}{$4\nu_1+2\nu_2+3\nu_4$}    &&&   \multicolumn{5}{c}{$2\nu_2+\nu_3+3\nu_4$}     &&&   \multicolumn{4}{c}{$\nu_3+5\nu_4$}            &&&   \multicolumn{1}{c}{$\nu_2+5\nu_4$}            &  \\
& \multicolumn{3}{c}{}                     &&&   \multicolumn{1}{c}{}                          &&&   \multicolumn{5}{c}{}                          &&&   \multicolumn{4}{c}{}                          &&&   \multicolumn{1}{c}{}                          &  \\
\hline
&&&&&&&&&&&&&&&&&&&&&&& \\
% A        B       C        D        E      F      G      H      I        J       K       L       M        N
&\1     & \2    & \3   &&& \4   &&& \5   & \6   & \7   & \8   & \9   &&& \la   & \lb   & \lc   & \ld  &&& \li   & \\
&&&&&&&&&&&&&&&&&&&&&&& \\
\hline
&&&&&&&&&&&&&&&&&&&&&&& \\
&\A     & \ab   & \ac  &&& \ad  &&& \ai  & \af  & \ag  & \ah  & \aii &&& \aj   & \ak   & \al   & \am  &&& \an   & \\
&\AR    &       &      &&&      &&&      &      &      &      &      &&&       &       &       &      &&&       & \\
&\AAA   &       &      &&&      &&&      &      &      &      &      &&&       &       &       &      &&&       & \\
&&&&&&&&&&&&&&&&&&&&&&& \\
\hline
&&&&&&&&&&&&&&&&&&&&&&& \\
&       & \B    & \bc  &&& \bd  &&& \be  & \br  & \bg  & \bh  & \bi  &&& \bj   & \bk   & \bl   & \bm  &&& \bn   & \\
&       & \BB   &      &&&      &&&      &      &      &      &      &&&       &       &       &      &&&       & \\
&       & \BBB  &      &&&      &&&      &      &      &      &      &&&       &       &       &      &&&       & \\
&&&&&&&&&&&&&&&&&&&&&&& \\
\hline
&&&&&&&&&&&&&&&&&&&&&&& \\
&       &       & \C   &&& \cd  &&& \ce  & \cf  & \cg  & \ch  & \ci  &&& \cj   & \ck   & \cl   & \cm  &&& \cn   & \\
&       &       & \CC  &&&      &&&      &      &      &      &      &&&       &       &       &      &&&       & \\
&&&&&&&&&&&&&&&&&&&&&&& \\
\hline
&&&&&&&&&&&&&&&&&&&&&&& \\
&       &       &      &&& \D   &&& \de  & \df  & \dg  & \dr  & \di  &&& \djr  & \dk   & \dl   & \dm  &&& \dn   & \\
&&&&&&&&&&&&&&&&&&&&&&& \\
\hline
&&&&&&&&&&&&&&&&&&&&&&& \\
&       &       &      &&&      &&& \E   & \ef  & \eg  & \eh  & \ei  &&& \ej  & \ek   & \el   & \er  &&& \en   & \\
&       &       &      &&&      &&& \EE  &      &      &      &      &&&      &       &       &      &&&       & \\
&       &       &      &&&      &&& \EEE & \eff &      &      & \eii &&&      &       &       &      &&&       & \\
&       &       &      &&&      &&& \EEEE&      &      &      &      &&&      &       &       &      &&&       & \\
&&&&&&&&&&&&&&&&&&&&&&& \\
\hline
&&&&&&&&&&&&&&&&&&&&&&& \\
&       &       &      &&&      &&&      & \F   & \fg  & \fh  & \fr  &&& \fj  & \fk   & \fl   & \fm  &&& \fn   & \\
&       &       &      &&&      &&&      & \FF  &      &      &      &&&      &       &       &      &&&       & \\
&       &       &      &&&      &&&      & \FFF &      &      &      &&&      &       &       &      &&&       & \\
&&&&&&&&&&&&&&&&&&&&&&& \\
\hline
&&&&&&&&&&&&&&&&&&&&&&& \\
&       &       &      &&&      &&&      &      & \G   & \gh  & \gi  &&& \gj  & \gk   & \gl   & \gm  &&& \gn   & \\
&       &       &      &&&      &&&      &      & \GG  &      &      &&&      &       &       &      &&&       & \\
&       &       &      &&&      &&&      &      & \GGG &      &      &&&      &       &       &      &&&       & \\
&&&&&&&&&&&&&&&&&&&&&&& \\
\hline
&&&&&&&&&&&&&&&&&&&&&&& \\
&       &       &      &&&      &&&      &      &      & \HH  & \hi  &&& \hj  & \hk   & \hl   & \hm  &&& \hn   & \\
&       &       &      &&&      &&&      &      &      & \HHH &      &&&      &       &       &      &&&       & \\
&&&&&&&&&&&&&&&&&&&&&&& \\
\hline
&&&&&&&&&&&&&&&&&&&&&&& \\
&       &       &      &&&      &&&      &      &      &      & \I   &&& \ir  & \ik   & \il   & \im  &&& \io   & \\
&       &       &      &&&      &&&      &      &      &      & \II  &&&      &       &       &      &&&       & \\
&       &       &      &&&      &&&      &      &      &      & \III &&&      &       &       &      &&&       & \\
&       &       &      &&&      &&&      &      &      &      & \IIII&&&      &       &       &      &&&       & \\
&&&&&&&&&&&&&&&&&&&&&&& \\
\hline
&&&&&&&&&&&&&&&&&&&&&&& \\
&       &       &      &&&      &&&      &      &      &      &      &&& \J   & \jk   & \jl   & \jm  &&& \jn   & \\
&       &       &      &&&      &&&      &      &      &      &      &&& \JJ  &       &       &      &&&       & \\
&       &       &      &&&      &&&      &      &      &      &      &&& \JJJ & \jkk  &       &      &&&       & \\
&&&&&&&&&&&&&&&&&&&&&&& \\
\hline
&&&&&&&&&&&&&&&&&&&&&&& \\
&       &       &      &&&      &&&      &      &      &      &      &&&      & \K    & \kl   & \km  &&& \kn   & \\
&       &       &      &&&      &&&      &      &      &      &      &&&      & \KK   &       &      &&&       & \\
&       &       &      &&&      &&&      &      &      &      &      &&&      & \KKK  &       &      &&&       & \\
&&&&&&&&&&&&&&&&&&&&&&& \\
\hline
&&&&&&&&&&&&&&&&&&&&&&& \\
&       &       &      &&&      &&&      &      &      &      &      &&&      &       & \LL   & \lm  &&& \lr   & \\
&       &       &      &&&      &&&      &      &      &      &      &&&      &       & \LLL  &      &&&       & \\
&       &       &      &&&      &&&      &      &      &      &      &&&      &       & \LLLL &      &&&       & \\
&&&&&&&&&&&&&&&&&&&&&&& \\
\hline
&&&&&&&&&&&&&&&&&&&&&&& \\
&       &       &      &&&      &&&      &      &      &      &      &&&      &       &       & \M   &&& \mn   & \\
&       &       &      &&&      &&&      &      &      &      &      &&&      &       &       & \MM  &&&       & \\
&&&&&&&&&&&&&&&&&&&&&&& \\
\hline
&&&&&&&&&&&&&&&&&&&&&&& \\
&       &       &      &&&      &&&      &      &      &      &      &&&      &       &       &      &&& \N    & \\
&&&&&&&&&&&&&&&&&&&&&&& \\
\hline
\end{tabular}
}%
}
\endgroup
\end{table}


% Floating table for part2; requires graphicx.
\begin{table}[p]
\centering
\caption{Тетраэдрическое расщепление и резонансные взаимодействия для набора полос $4\nu_3$, $\nu_1+3\nu_3$, $2\nu_1+2\nu_3$, $3\nu_1+\nu_3(F_2)$.}
\label{tab:si_hamiltonian_elements_f2}
\begingroup
  \def\A{$20G_{33}-\frac{156}{7}T_{33}$}
  \def\ab{$40\frac{\sqrt{6}}{7}T_{33}$}
  \def\ac{$-3\frac{\sqrt{35}}{140}F_{1333}$}
  \def\ad{$-3\frac{\sqrt{210}}{70}F_{1333}$}
  \def\ai{$0$}
  \def\af{$0$}
  \def\B{$6G_{33}-\frac{152}{7}T_{33}$}
  \def\bc{$\sqrt{\frac{6}{35}}F_{1333}$}
  \def\bd{$\frac{13\sqrt{35}}{140}F_{1333}$}
  \def\be{$-\frac{\sqrt{7}}{2}F_{1333}$}
  \def\br{$0$}
  \def\C{$12G_{33}+12T_{33}$}
  \def\cd{$-8\sqrt{6}T_{33}$}
  \def\ce{$-\frac{\sqrt{30}}{10}F_{1333}$}
  \def\cf{$0$}
  \def\D{$2G_{33}$}
  \def\de{$-\frac{1}{\sqrt{5}}F_{1333}$}
  \def\df{$-\frac{\sqrt{15}}{2}F_{1333}$}
  \def\E{$6G_{33}-8T_{33}$}
  \def\ef{$\frac{\sqrt{3}}{4}F_{1333}$}
  \def\F{$2G_{33}$}
  \def\1{$l_3=4$ }
  \def\2{$l_3=2$ }
  \def\3{$l_3=3$ }
  \def\4{$l_3=1$ }
  \def\5{$l_3=2$ }
  \def\6{$l_3=1$ }
\scriptsize
\rotatebox{90}{%
\resizebox{0.80\textheight}{!}{%
\begin{tabular}{|r c|c l|r c|c l|r c l|r c l|}\hline
& \multicolumn{2}{c}{}                     &&&   \multicolumn{2}{c}{}                           &&&   \multicolumn{1}{c}{}                           &&&   \multicolumn{1}{c}{}                           &  \\
& \multicolumn{2}{c}{$4\nu_3(F_2)$}        &&&   \multicolumn{2}{c}{$\nu_1+3\nu_3(F_2)$}        &&&   \multicolumn{1}{c}{$2\nu_1+2\nu_3(F_2)$}       &&&   \multicolumn{1}{c}{$3\nu_1+\nu_3(F_2)$}        &  \\
& \multicolumn{2}{c}{}                     &&&   \multicolumn{2}{c}{}                           &&&   \multicolumn{1}{c}{}                           &&&   \multicolumn{1}{c}{}                           &  \\
\hline
&&&&&&&&&&&&& \\
% A       B          C      D        E        F      
&\1     & \2    &&& \3   & \4   &&& \5   &&& \6   & \\
&&&&&&&&&&&&& \\
&&&&&&&&&&&&& \\
\hline
&&&&&&&&&&&&& \\
&\A     & \ab   &&& \ac  & \ad  &&& \ai  &&& \af  & \\
&&&&&&&&&&&&& \\
\hline
&&&&&&&&&&&&& \\
&       & \B    &&& \bc  & \bd  &&& \be  &&& \br  & \\
&&&&&&&&&&&&& \\
\hline
&&&&&&&&&&&&& \\
&       &       &&& \C   & \cd  &&& \ce  &&& \cf  & \\
&&&&&&&&&&&&& \\
\hline
&&&&&&&&&&&&& \\
&       &       &&&      & \D   &&& \de  &&& \df  & \\
&&&&&&&&&&&&& \\
\hline
&&&&&&&&&&&&& \\
&       &       &&&      &      &&& \E   &&& \ef  & \\
&&&&&&&&&&&&& \\
\hline
&&&&&&&&&&&&& \\
&       &       &&&      &      &&&      &&& \F   & \\
&&&&&&&&&&&&& \\
\hline
\end{tabular}
}%
}
\endgroup
\end{table}


\subsection{Учет эффекта локальных мод}

Молекула SiH$_4$ относится к классу молекул, для которых приближение локальных мод является
существенным \cite{301,302,303,304,305}. В предельном случае это означает, что масса центрального атома
значительно превышает массу периферических атомов, а взаимодействие между различными растягивающими
координатами и между растягивающими и деформационными координатами мало по сравнению с собственно
растягивающими силовыми постоянными. В таком приближении для растягивающих частот выполняются
соотношения
\begin{equation}
\omega_1 = \omega + 3\lambda, \qquad \omega_3 = \omega - \lambda,
\label{eq:si_local_frequencies}
\end{equation}
где $|\lambda/\omega|\ll 1$. Кроме того, параметры, отвечающие за ангармонизм и локально-модовое
расщепление, связаны приближенными соотношениями
\begin{equation}
x_{11} \simeq \frac{1}{4}x_{13} \simeq \frac{5}{9}x_{33} \simeq
-\frac{5}{3}G_{33} \simeq -5T_{33} \simeq \frac{1}{4}d_{1133} \simeq
\frac{1}{16}d_{1333}.
\label{eq:si_local_parameters}
\end{equation}
В реальной молекуле эти равенства не являются строгими, поскольку растягивающие и деформационные
колебания взаимодействуют. Тем не менее анализ показывает, что для ${}^{28}$SiH$_4$ суммы
$x_{12}+3x_{14}/2$ и $x_{23}+3x_{34}/2$ отличаются достаточно мало, чтобы локально-модовая интерпретация
оставалась применимой. По результатам подгонки для ${}^{28}$SiH$_4$ получены значения
$x_{12}+3x_{14}/2=-13.055$ см$^{-1}$ и $x_{23}+3x_{34}/2=-15.390$ см$^{-1}$. Отсюда следует, что параметр
$\lambda$ для SiH$_4$ больше, чем для GeH$_4$, однако остается малым по сравнению с основной
растягивающей частотой. Это подтверждает возможность использовать локально-модовые соотношения как
физическое ограничение при выборе набора спектроскопических параметров.

\subsection{Определение параметров для ${}^{28}$SiH$_4$}

Для основного изотополога ${}^{28}$SiH$_4$ в качестве исходной информации использовались 49
экспериментальных центров полос, опубликованных в работах \cite{4302,4301,4303,4300,5312,4312,zhu1,zhu2,zhu3},
а также 23 центра полос из базы STDS, использованные с пониженным весом \cite{dijon}. Эти данные
позволили определить набор колебательных спектроскопических параметров, включающий гармонические
частоты, ангармонические коэффициенты, параметры тетраэдрического расщепления и параметры
резонансных взаимодействий. Взвешенная подгонка выполнена методом наименьших квадратов с учетом
аналитических матричных элементов из \cite{ulen1,ulen2}.

Полученный набор параметров воспроизводит исходные центры полос ${}^{28}$SiH$_4$ со среднеквадратичным
отклонением $d_{\textrm{rms}}=0.28$ см$^{-1}$. Для проверки качества решения результаты были сопоставлены с
расчетом \cite{wang}, основанным на силовом поле \cite{lee_1}. Для этого набора параметров отклонение
составляет $d_{\textrm{rms}}=2.21$ см$^{-1}$, то есть более чем в восемь раз превышает значение, полученное в
настоящем анализе. Такое сравнение показывает, что полуэмпирическое решение, основанное на
экспериментальных центрах полос, лучше описывает известную колебательную структуру молекулы SiH$_4$ и
должно надежнее предсказывать положения более высоких полос.

Дополнительным критерием физической корректности параметров служит выполнение локально-модовых
соотношений. В найденном наборе параметры $d_{1133}$ и $d_{1333}$ согласуются с ожидаемыми оценками
$d_{1133}\simeq4x_{11}$ и $d_{1333}\simeq16x_{11}$, тогда как в альтернативном наборе \cite{wang} эти
соотношения нарушаются. Это важно не только для воспроизведения уже известных центров полос, но и для
экстраполяции энергетической структуры в область, где прямые экспериментальные данные пока отсутствуют.

\input{tables/si_table_band_centers_fit.tex}

% Reworked from Tab_4.tex as non-floating longtable.
\begingroup
\small
\setlength{\tabcolsep}{4pt}
\renewcommand{\arraystretch}{1.05}
\begin{longtable}{lcccc}
\caption{Колебательные спектроскопические параметры изотопологов силана, полученные из решения обратной задачи. Значения приведены в см$^{-1}$.}\label{tab:si_vibrational_parameters}\\
\hline
Параметр & ${}^{28}$SiH$_4$ (настоящая работа) & ${}^{28}$SiH$_4$ (\cite{wang}$^{b)}$) & ${}^{29}$SiH$_4$ & ${}^{30}$SiH$_4$ \\
\hline
\endfirsthead
\hline
\multicolumn{5}{c}{\small\slshape Продолжение таблицы~\ref{tab:si_vibrational_parameters}}\\
\hline
Параметр & ${}^{28}$SiH$_4$ (настоящая работа) & ${}^{28}$SiH$_4$ (\cite{wang}$^{b)}$) & ${}^{29}$SiH$_4$ & ${}^{30}$SiH$_4$ \\
\hline
\endhead
\hline
\multicolumn{5}{r}{\small\slshape Продолжение следует}\\
\endfoot
\hline\hline
\endlastfoot
$x_{11}$ & -8.56(15) & -8.364 & -8.53 & -8.50(18) \\n$x_{12}$ & -5.45(37) & -6.157 & -5.45 & -5.45 \\n$x_{13}$ & -33.27(15) & -32.803 & -33.47 & -33.67(26) \\n$x_{14}$ & -5.07(37) & -4.972 & -5.07 & -5.07 \\n$x_{22}$ & -1.115(51) & -0.847 & -1.102 & -1.089(35) \\n$x_{23}$ & -6.57(26) & -7.857 & -6.57 & -6.57 \\n$x_{24}$ & 0.148(49) & 0.637 & 0.148 & 0.148 \\n$x_{33}$ & -14.949(90) & -14.935 & -15.001 & -15.053(39) \\n$x_{34}$ & -5.88(22) & -6.188 & -5.88 & -5.88 \\n$x_{44}$ & -2.833(36) & -2.130 & -2.760 & -2.687(28) \\n$G_{22}$ & 1.549(41) & 1.351 & 1.544 & 1.539(41) \\n$G_{33}$ & 5.438(48) & 5.452 & 5.448 & 5.458(32) \\n$G_{34}$ & 0.20(11) & 0.605 & 0.20 & 0.20 \\n$G_{44}$ & 2.132(25) & 1.987 & 2.117 & 2.102(37) \\n$T_{23}$ & 0.053(24) & 0.007 & 0.053 & 0.053 \\n$T_{24}$ & -0.2933(50) & -0.292 & -0.3122 & -0.3311(48) \\n$T_{33}$ & 1.755(13) & 1.745 & 1.750 & 1.745(29) \\n$T_{34}$ & -0.217(39) & -0.138 & -0.217 & -0.217 \\n$T_{44}$ & 0.2024(39) & 0.094 & 0.1882 & 0.1740(43) \\n$S_{34}$ & 0.488(42) & 0.322 & 0.488 & 0.488 \\n$d_{1133}$ & -33.90(19) & -8.428 & -33.93 & -33.96(36) \\n$d_{1333}$ & -137.12(38) & --- & -137.54 & -137.96(48) \\n$d_{2244s}$ & 8.75(55) & -0.998 & 8.75 & 8.75 \\n$d_{2244t}$ & 2.9(11) & -0.340 & 2.9 & 2.9 \\n$d_{2444 }$ & --- & 1.743 & --- & --- \\n$c_{122 } $ & --- & 6.021 & --- & --- \\n$c_{144 } $ & --- & 10.768 & --- & --- \\n$c_{344 } $ & --- & -17.502 & --- & --- \\n$c_{234 } $ & --- & 8.397 & --- & --- \\n\end{longtable}
\vspace{0.5ex}
\begin{minipage}{0.95\textwidth}
\scriptsize
$^{a)}$ Обозначения параметров соответствуют работе \cite{hecht}.\quad
$^{b)}$ Воспроизведено из работы \cite{wang}.
\end{minipage}
\endgroup


На основе найденных параметров был рассчитан полный набор центров колебательных полос ${}^{28}$SiH$_4$
до 9000 см$^{-1}$ для всех состояний с полиадным числом $N\leq4$. В итоговый список вошло 2125 центров
полос для данного изотополога. Поскольку полный список достаточно велик, в тексте диссертации приведен
фрагмент такого списка, тогда как полный набор целесообразно вынести в приложение.

\subsection{Изотопическое замещение и анализ ${}^{29}$SiH$_4$ и ${}^{30}$SiH$_4$}

Для изотопологов ${}^{29}$SiH$_4$ и ${}^{30}$SiH$_4$ количество доступных экспериментальных центров полос
меньше, чем для ${}^{28}$SiH$_4$. Поэтому при их анализе дополнительно использовались соотношения теории
изотопозамещения \cite{byk}. Пусть $\omega_i$ --- гармонические частоты материнской молекулы, а
$\omega'_k$ --- частоты изотопически замещенной молекулы. Тогда связь между ними может быть записана в
виде
\begin{equation}
\sum_i \omega_i^2 \alpha_i^k = \alpha_j^k {\omega'_k}^2,
\label{eq:si_isotope_frequencies}
\end{equation}
где коэффициенты $\alpha_i^k$ определяются из условий преобразования нормальных координат. Матрица,
связывающая нормальные координаты двух изотопологов, удовлетворяет соотношению
\begin{equation}
\sum_k \alpha_i^k \alpha_j^k = A_{ij},
\label{eq:si_alpha_matrix}
\end{equation}
а элементы $A_{ij}$ выражаются через массы ядер до и после замещения и через коэффициенты
преобразования нормальных координат:
\begin{equation}
A_{ij}=\delta_{ij}-\sum_{N\beta}\frac{m'_N-m_N}{m'_N}l_{N\beta i}l_{N\beta j}.
\label{eq:si_aij}
\end{equation}
В случае силана замещается центральный атом кремния, поэтому влияние изотопозамещения на параметры
оказывается достаточно регулярным. Анализ этих соотношений приводит к простому приближенному правилу:
для одноименных параметров $P_n$ трех изотопологов выполняется
\begin{equation}
{}^{29}P_n \simeq \frac{1}{2}\left({}^{28}P_n+{}^{30}P_n\right).
\label{eq:si_isotope_parameter_relation}
\end{equation}
Это соотношение использовалось как дополнительное условие при совместной подгонке параметров
${}^{29}$SiH$_4$ и ${}^{30}$SiH$_4$.

В расчет были включены 32 экспериментальных центра полос ${}^{29}$SiH$_4$ и ${}^{30}$SiH$_4$, полученные из
работ \cite{4302,4301,4303,4300}. Так как для этих изотопологов отсутствуют экспериментальные данные о
некоторых растяжительно-деформационных комбинационных полосах, часть параметров была зафиксирована на
значениях, найденных для ${}^{28}$SiH$_4$. К таким параметрам относятся $x_{12}$, $x_{14}$, $x_{23}$,
$x_{34}$, $G_{34}$, $T_{23}$, $T_{34}$ и $S_{34}$. Кроме того, на основании анализа статистических
неопределенностей были дополнительно зафиксированы параметры $x_{14}$, $d_{1333}$, $d_{2244s}$ и
$d_{2244t}$. Остальные параметры варьировались в совместной процедуре наименьших квадратов.

В результате был получен набор из 18 варьируемых параметров, который воспроизводит 32 исходных центра
полос двух изотопологов со среднеквадратичным отклонением $d_{\textrm{rms}}=0.06$ см$^{-1}$. Такая точность
показывает, что совместное использование экспериментальных данных и изотопических ограничений
позволяет устойчиво решить обратную задачу даже при ограниченном числе исходных полос.

% Reworked from tab_5.tex as non-floating longtable.
\begingroup
\small
\setlength{\tabcolsep}{5pt}
\renewcommand{\arraystretch}{1.05}
\begin{longtable}{lccc}
\caption{Фрагмент рассчитанного списка колебательных уровней молекулы ${}^{M}$SiH$_4$ ($M=28,29,30$). Значения приведены в см$^{-1}$.}\label{tab:si_predicted_band_centers}\\
\hline
Состояние & $M=28$ & $M=29$ & $M=30$ \\
\hline
\endfirsthead
\hline
\multicolumn{4}{c}{\small\slshape Продолжение таблицы~\ref{tab:si_predicted_band_centers}}\\
\hline
Состояние & $M=28$ & $M=29$ & $M=30$ \\
\hline
\endhead
\hline
\multicolumn{4}{r}{\small\slshape Продолжение следует}\\
\endfoot
\hline\hline
\endlastfoot
$( 0 1 1 2;~   1 1 2;~   3 F_2  )$&  4965.32 &  4961.24 &  4957.23 \\
$( 0 1 1 2;~   1 1 2;~   4 F_2  )$&  4967.88 &  4964.25 &  4960.63 \\
$( 1 1 0 2;~   1 0 2;~     F_2  )$&  4968.76 &  4966.56 &  4964.35 \\
$( 0 1 1 2;~   1 1 2;~   5 F_2  )$&  4971.82 &  4968.02 &  4964.24 \\
$( 0 2 1 1;~   0 1 1;~   1 F_2  )$&  5017.93 &  5015.04 &  5012.15 \\
$( 1 2 0 1;~   0 0 1;~   1 F_2  )$&  5018.63 &  5017.13 &  5015.62 \\
$( 0 2 1 1;~   2 1 1;~   3 F_2  )$&  5027.06 &  5024.47 &  5021.88 \\
$( 0 2 1 1;~   2 1 1;~   2 F_2  )$&  5029.32 &  5026.95 &  5024.59 \\
$( 1 2 0 1;~   2 0 1;~   2 F_2  )$&  5029.88 &  5028.87 &  5027.87 \\
$( 0 3 1 0;~   1 1 0;~   2 F_2  )$&  5075.02 &  5073.73 &  5072.45 \\
$( 0 3 1 0;~   3 1 0;~   1 F_2  )$&  5084.80 &  5083.51 &  5082.23 \\
$( 0 0 2 1;~   0 0 1;~   1 F_2  )$&  5210.28 &  5206.79 &  5203.25 \\
$( 1 0 1 1;~   0 1 1;~     F_2  )$&  5213.35 &  5209.79 &  5275.12 \\
$( 1 1 1 0;~   1 1 0;~     F_2  )$&  5267.84 &  5334.32 &  5332.26 \\
$( 2 0 0 1;~   0 0 1;~     F_2  )$&  5277.25 &  5275.20 &  5272.98 \\
$( 0 0 2 1;~   0 2 1;~   2 F_2  )$&  5281.50 &  5278.64 &  5206.22 \\
$( 0 0 2 1;~   0 2 1;~   3 F_2  )$&  5282.51 &  5277.86 &  5274.08 \\
$( 0 1 2 0;~   1 2 0;~     F_2  )$&  5336.40 &  5265.55 &  5263.25 \\
$( 0 0 0 6;~   0 0 0;~   3 A_1  )$&  5360.42 &  5355.67 &  5350.90 \\
$( 0 0 0 6;~   0 0 2;~   1 F_2  )$&  5369.09 &  5364.46 &  5359.82 \\
$( 0 3 0 4;~   3 0 4;~   2 F_2  )$&  6559.59 &  6555.61 &  6551.64 \\
$( 0 1 0 6;~   1 0 6;~   6 F_1  )$&  6563.66 &  6556.34 &  6548.98 \\
$( 0 1 0 6;~   1 0 6;~   1 A_1  )$&  6564.47 &  6556.06 &  6547.69 \\
$( 0 4 0 3;~   0 0 1;~   1 F_2  )$&  6564.96 &  6560.19 &  6555.36 \\
$( 3 0 0 0;~   0 0 0;~   1 A_1  )$&  6565.20 &  6562.95 &  6560.79 \\
$( 3 0 0 0;~   0 0 0;~   1 A_1  )$&  6565.20 &  6562.95 &  6560.79 \\
$( 0 4 0 3;~   0 0 1;~   1 F_2  )$&  6565.67 &  6561.11 &  6556.49 \\
$( 0 1 0 6;~   1 0 6;~   6  E   )$&  6566.04 &  6558.32 &  6550.65 \\
$( 0 1 0 6;~   1 0 6;~   5 F_1  )$&  6568.38 &  6561.16 &  6554.07 \\
$( 0 2 0 5;~   2 0 3;~   2  E   )$&  6568.94 &  6563.68 &  6558.42 \\
$( 0 3 0 4;~   3 0 4;~   3 F_2  )$&  6570.40 &  6566.53 &  6562.69 \\
$( 0 0 3 0;~   0 3 0;~   1 F_2  )$&  6570.41 &  6566.68 &  6562.98 \\
$( 0 0 3 0;~   0 3 0;~   1 F_2  )$&  6570.41 &  6566.68 &  6562.98 \\
$( 0 3 0 4;~   3 0 4;~   3 F_2  )$&  6570.44 &  6566.56 &  6562.73 \\
$( 0 2 0 5;~   0 0 3;~   1 F_1  )$&  6572.42 &  6566.26 &  6560.04 \\
$( 0 2 0 5;~   2 0 3;~   4 F_1  )$&  6576.46 &  6570.42 &  6564.37 \\
$( 0 1 0 6;~   1 0 6;~   7  E   )$&  6578.07 &  6570.76 &  6563.46 \\
$( 0 1 0 6;~   1 0 6;~   3 A_2  )$&  6581.77 &  6574.29 &  6566.82 \\
$( 0 4 0 3;~   2 0 1;~   3 F_2  )$&  6582.13 &  6578.20 &  6574.22 \\
$( 0 4 0 3;~   2 0 1;~   3 F_2  )$&  6584.48 &  6580.92 &  6577.29 \\
$( 0 2 0 5;~   2 0 3;~   5 F_1  )$&  6587.33 &  6581.24 &  6575.19 \\
$( 0 4 0 3;~   2 0 3;~   4 F_2  )$&  6595.83 &  6592.30 &  6588.80 \\
$( 0 4 0 3;~   0 0 3;~   2 F_2  )$&  6597.19 &  6594.06 &  6633.96 \\
$( 0 4 0 3;~   4 0 1;~   6 F_2  )$&  6604.98 &  6601.76 &  6598.55 \\
\end{longtable}
\endgroup


Итогом анализа является согласованный набор спектроскопических параметров для трех стабильных
изотопологов силана. Эти параметры позволяют рассчитать 6375 центров колебательных полос различных
симметрий, то есть по 2125 центров для каждого изотополога. Полученный список может использоваться как
основа для дальнейшего моделирования высокоразрешенных спектров SiH$_4$, для уточнения параметров
внутримолекулярной поверхности потенциальной энергии и для последующего построения более полных
спектроскопических баз данных.




\section{Параметры ангармоничности, тэтраэдрического расщипления и резонансов}\label{sec:ch2/params}

\section{Потенциальная функция молекул GeH$_4$ и SiH$_4$}\label{sec:ch2/pot-func}
